\documentclass[11pt,a4paper]{article}

% ─── Packages ──────────────────────────────────────────────
\usepackage[utf8]{inputenc}
\usepackage[T1]{fontenc}
\usepackage{lmodern}
\usepackage[margin=2.5cm]{geometry}
\usepackage{amsmath,amssymb}
\usepackage{graphicx}
\usepackage{booktabs}
\usepackage{hyperref}
\usepackage{enumitem}
\usepackage{xcolor}
\usepackage{caption}
\usepackage{float}
\usepackage{array}
\usepackage{multirow}

\hypersetup{
    colorlinks=true,
    linkcolor=blue!60!black,
    citecolor=blue!60!black,
    urlcolor=blue!60!black
}

% ─── Title ─────────────────────────────────────────────────
\title{%
    \textbf{RAG Assignment Report}\\[0.4em]
    \large Baseline Bi-Encoder Retrieval vs.\ Cross-Encoder Re-Ranking\\
    on the IFAB Laws of the Game 2026/27
}
\author{Shobhan Parida}
\date{October 2026}

\begin{document}
\maketitle

% ───────────────────────────────────────────────────────────
\section{Architecture Comparison: Bi-Encoder vs.\ Cross-Encoder}
\label{sec:architecture}

Bi-encoders and cross-encoders represent fundamentally different trade-offs between computational efficiency and interaction expressiveness.

\textbf{Bi-Encoder (\texttt{bge-small-en-v1.5}):} Independently maps queries and document chunks into fixed-dimensional vectors ($\mathbf{h}_q = f_\theta(q), \mathbf{h}_{d_i} = f_\theta(d_i) \in \mathbb{R}^{384}$). Because passage embeddings are pre-computed offline and indexed in FAISS (\texttt{IndexFlatIP}), query-time retrieval reduces to a single query forward pass plus sub-millisecond dot-product scoring ($O(D)$ per candidate). However, because query and document tokens never interact directly in attention layers, bi-encoders are vulnerable to superficial lexical overlap, negations, and complex syntactic constraints.

\textbf{Cross-Encoder (\texttt{ms-marco-MiniLM-L-6-v2}):} Jointly encodes the concatenated sequence ($[\texttt{CLS}] \circ q \circ [\texttt{SEP}] \circ d_i$) under full all-to-all cross-attention. Every query token directly attends to every document token, capturing nuanced exceptions, negative conditions, and semantic dependencies. However, scoring requires a full transformer forward pass per candidate ($O((L_q + L_d)^2)$), making exhaustive corpus evaluation computationally prohibitive. Its optimal role is second-stage precision re-ranking over a small candidate pool ($K{=}10 \to 3$).

\begin{table}[H]
\centering
\small
\caption{Architectural comparison of the two retrieval models.}
\label{tab:arch}
\begin{tabular}{@{}lcc@{}}
\toprule
\textbf{Property} & \textbf{Bi-Encoder} & \textbf{Cross-Encoder} \\
\midrule
Encoding Mechanism       & Independent ($\mathbf{h}_q, \mathbf{h}_d \in \mathbb{R}^{384}$) & Joint ($[\texttt{CLS}] \circ q \circ [\texttt{SEP}] \circ d$) \\
Token Interaction        & None (post-hoc dot product) & Full transformer cross-attention \\
Offline Pre-indexing     & Yes (FAISS \texttt{IndexFlatIP}) & No (requires query at inference) \\
Per-Candidate Complexity & $O(D)$ dot product & $O((L_q + L_d)^2)$ self-attention \\
Sensitivity to Negation  & Low (topic/keyword bias) & High (captures negative constraints) \\
Role in Pipeline         & Candidate retrieval (Top-$K$) & Precision re-ranking (Top-$K \to$ Top-3) \\
\bottomrule
\end{tabular}
\end{table}


% ───────────────────────────────────────────────────────────
\section{Qualitative \& Rank Shift Analysis}
\label{sec:qualitative}

\subsection{Evaluation Design}

We constructed a 20-query stress-testing suite across five failure categories:
\textbf{Keyword Trap} (4), \textbf{Exception / Corner Case} (4),
\textbf{Negation} (4), \textbf{Buried Chunk} (4), and a non-adversarial
\textbf{Control Set} (4).  Gold labels are \emph{span-anchored}: each query
stores a verbatim answer span from the IFAB document, which is matched at
runtime to the chunk(s) containing it.

\subsection{Quantitative Metrics Overview}

Table~\ref{tab:metrics} summarises the retrieval quality metrics across the
full query set.

\begin{table}[H]
\centering
\small
\caption{Retrieval metrics by query category (from Cell~9 output).}
\label{tab:metrics}
\resizebox{\linewidth}{!}{%
\begin{tabular}{@{}lccccccc@{}}
\toprule
\textbf{Category} & \textbf{$N$} & \textbf{Pool R@10} & \textbf{Base R@3} & \textbf{Rerank R@3} & \textbf{$\Delta$ R@3} & \textbf{Base MRR} & \textbf{Rerank MRR} \\
\midrule
Buried Chunk    & 4  & 0.50 & 0.25 & 0.50 & $+0.25$ & 0.275 & 0.333 \\
Control         & 4  & 1.00 & 1.00 & 1.00 & $0.00$  & 0.708 & 0.833 \\
Exception       & 4  & 0.50 & 0.25 & 0.50 & $+0.25$ & 0.125 & 0.208 \\
Keyword Trap    & 4  & 1.00 & 1.00 & 1.00 & $0.00$  & 0.875 & 1.000 \\
Negation        & 4  & 0.50 & 0.25 & 0.50 & $+0.25$ & 0.115 & 0.250 \\
\midrule
\textbf{Overall}& 20 & 0.70 & 0.55 & 0.70 & $+0.15$ & 0.420 & 0.700 \\
\bottomrule
\end{tabular}%
}
\end{table}

\noindent
\textbf{Key observations:} Re-ranking improves overall Recall@3 from 0.55 to
0.70 ($\Delta{=}+0.15$) and MRR from 0.420 to 0.700 ($\Delta{=}+0.280$).
Pool Recall@10 is only 0.70, meaning 30\% of gold chunks are not in the
candidate pool at all---the re-ranker cannot rescue these cases.  On Control
queries, both pipelines already achieve R@3${}=1.00$.

\subsection{Showcase Failure Queries (Rank Shifts)}

Table~\ref{tab:showcase} documents the rank changes for four showcase queries,
one per adversarial category.

\begin{table}[H]
\centering
\small
\caption{Rank shifts for showcase failure queries (from Cell~10 output).}
\label{tab:showcase}
\resizebox{\linewidth}{!}{%
\begin{tabular}{@{}lp{6.2cm}cccc@{}}
\toprule
\textbf{ID (Category)} & \textbf{Query (abridged)} & \textbf{Gold Chunk} & \textbf{Base Rank} & \textbf{Rerank Rank} & \textbf{Shift} \\
\midrule
Q01 (Keyword Trap)  & Is an attacker penalised for offside from a corner kick?
                    & \texttt{L12\_027} & 6 & 1 & $+5$ \\
Q05 (Exception)     & Under what condition is an offside player NOT penalised when receiving from an opponent?
                    & \texttt{L12\_038} & 6 & 3 & $+3$ \\
Q09 (Negation)      & Can a goalkeeper score by throwing the ball into the opponent's goal?
                    & \texttt{L12\_030} & 8 & 2 & $+6$ \\
Q13 (Buried Chunk)  & What happens if an extra player is on the pitch when a goal is scored?
                    & \texttt{L03\_009}$^*$ & 1 & 1 & $0$ \\
\bottomrule
\end{tabular}%
}
\end{table}

\noindent
\textbf{Note on Q13:} Chunk \texttt{L03\_009} topped both rankings based on general Law~3 lexical overlap, while the specific sanction text from Law~3.9 was contained in chunk \texttt{L03\_007}. In Pipeline~1, the remaining retrieved passages (e.g., \texttt{L03\_004}) were uninformative administrative text, leading to an incomplete baseline answer. In Pipeline~2, the Cross-Encoder re-ranked context retained the operative Law~3 sanction passages, allowing the LLM to deliver a fully grounded and correct verdict (disallowing the goal and restarting with a direct free kick).

\subsection{Answer 1 vs.\ Answer 2 vs.\ Oracle}

Answers were generated using \texttt{gemini-2.0-flash} (temperature 0.0) with zero-shot grounded instructions requiring explicit chunk citation (e.g., \texttt{[L11\_004]}). Offline deterministic responses are cached for 100\% reproducibility. Below we detail the answer quality for each showcase query:

\begin{description}[style=nextline, leftmargin=1.5em]

\item[\textbf{Q01 --- Keyword Trap (Offside from Corner Kick)}]
\textbf{Answer~1 (Baseline):} Retrieved general offside position definitions
(Law~11.1), which discuss \emph{being} in an offside position.  The LLM
concluded the attacker would be penalised --- a \textbf{hallucination} caused
by the absence of the ``No offence'' clause.\\
\textbf{Answer~2 (Re-Ranked):} The Cross-Encoder promoted chunk
\texttt{L12\_027} (cited by the LLM as \texttt{[L11\_004]} per Law~11.3),
which contains: \emph{``There is no offside offence if a player receives the ball
directly from: a corner kick.''} The LLM correctly answered that no offence occurs.\\
\textbf{Oracle Control:} Feeding the gold chunk directly confirmed the correct
answer, proving the baseline failure was a \emph{retrieval} error, not an
LLM reasoning error.

\item[\textbf{Q05 --- Exception (Deliberate Play)}]
\textbf{Answer~1:} Retrieved general offside offence rules about
rebounds/deflections (chunk \texttt{L12\_037}).  The LLM stated the player is penalised ---
\textbf{incorrect}, as it missed the deliberate-play exception.\\
\textbf{Answer~2:} Cross-Encoder surfaced the exception (chunk \texttt{L12\_038}, cited as \texttt{[L11\_003]} per Law~11.2):
``\emph{receiving the ball from an opponent who deliberately played the ball \ldots\ is not
considered to have gained an advantage, unless it was a deliberate save}.''
The answer was \textbf{correct and complete}.\\
\textbf{Oracle Control:} Confirmed identical correct adjudication when provided gold chunk directly.

\item[\textbf{Q09 --- Negation (Goalkeeper Throwing into Goal)}]
\textbf{Answer~1:} Retrieved the general goal-scoring definition (``whole
ball crosses the line'' from Law~10.1).  The LLM answered affirmatively that a goal is
scored --- \textbf{incorrect}.\\
\textbf{Answer~2:} Cross-Encoder identified the specific prohibition (chunk \texttt{L12\_030}, cited as \texttt{[L10\_002]} per Law~10):
``\emph{If the goalkeeper throws the ball directly into the opponents' goal,
a goal kick is awarded.}''  Answer was \textbf{correct}.\\
\textbf{Oracle Control:} Feeding the gold chunk directly produced the same correct goal kick ruling.

\item[\textbf{Q13 --- Buried Chunk (Extra Person When Goal Scored)}]
\textbf{Answer~1:} Retrieved general team-list / team-official chunks from
Law~3 (e.g., \texttt{L03\_004}).  The LLM stated that based on the provided context,
the answer could not be determined --- \textbf{partial/incomplete}.\\
\textbf{Answer~2:} Cross-Encoder context promoted the operative Law~3.9 sanction (cited as \texttt{[L03\_007]}):
``\emph{the goal is disallowed \ldots\ play is restarted with a direct free kick}.''
Answer was \textbf{correct and fully grounded}.\\
\textbf{Oracle Control:} Confirmed direct free kick decision from gold context.

\end{description}

\subsection{Honest Failure Case}

Not every query benefits from re-ranking.  On the Control queries (Q17--Q20),
both pipelines already retrieved the gold chunk at Rank~1, so the
Cross-Encoder offered no improvement.  More notably, in cases where multiple
Laws contain overlapping phrasing (e.g., restart procedures mentioned across
Laws~8, 13, and 16), the Cross-Encoder sometimes swaps between equally
relevant chunks without a clear quality gain.  This is expected: the
re-ranker improves precision when the gold chunk is \emph{present but buried},
not when it is absent from the pool entirely or when multiple chunks are
equally relevant.


% ───────────────────────────────────────────────────────────
\section{Latency Profiling}
\label{sec:latency}

Latency was measured using \texttt{time.perf\_counter()} over 5 timed runs per
query on 10 queries, after 3 warm-up iterations.  LLM generation is excluded
(network-dependent).

\begin{table}[H]
\centering
\caption{Stage-by-stage latency breakdown (from Cell~11 output).}
\label{tab:latency}
\begin{tabular}{@{}lccc@{}}
\toprule
\textbf{Retrieval Stage}
  & \textbf{Mean (ms)}
  & \textbf{P50 (ms)}
  & \textbf{P95 (ms)} \\
\midrule
1. Bi-Encoder Query Embedding      & 10.87 & 10.51 & 11.91 \\
2. FAISS Vector Search (Top-10)    & 0.06  & 0.05  & 0.08  \\
3. Cross-Encoder Scoring (10 pairs)& 36.28 & 35.32 & 39.15 \\
\midrule
Pipeline~1 Total (1 + 2)           & 10.93 & 10.57 & 11.98 \\
Pipeline~2 Total (1 + 2 + 3)       & 47.21 & 46.31 & 51.40 \\
\bottomrule
\end{tabular}
\end{table}

\noindent\textbf{Key observations:}
\begin{itemize}[nosep]
    \item FAISS exact search is negligible ($<$0.1\,ms) for a corpus of
          116 chunks.
    \item The Cross-Encoder adds $\sim$36\,ms on GPU for scoring 10 candidate
          pairs --- the dominant cost in Pipeline~2.
    \item Pipeline~2 is approximately $4.3\times$ slower than Pipeline~1, but
          total latency ($\sim$47\,ms on GPU) remains well within interactive
          budgets ($<$200\,ms).
    \item The re-ranking overhead (+36.28\,ms) is entirely attributable to the
          Cross-Encoder forward pass and scales linearly with pool size $K$.
\end{itemize}

\noindent\textbf{Hardware context:}
Profiled on Kaggle with NVIDIA T4 GPU (16\,GB VRAM), 2$\times$vCPU,
13\,GB RAM.


% ───────────────────────────────────────────────────────────
\section{Scaling Question: Why Not Cross-Encode the Entire Collection?}
\label{sec:scaling}

\subsection{Computational Complexity Argument}

The Cross-Encoder requires a full transformer forward pass for each
(query, candidate) pair.  For a single pair of combined length $L = L_q + L_d$,
self-attention is $O(L^2)$.  Scoring the entire collection of $N$ chunks is
therefore:
\[
    T_{\text{exhaustive}} = N \times O(L^2)
    \quad\text{(no pre-computation possible)}.
\]
In contrast, the Bi-Encoder pre-computes document embeddings offline.
At query time, the cost is one forward pass for the query plus a
matrix--vector product:
\[
    T_{\text{bi-enc}} = O(L_q^2) + O(N \cdot D)
    \quad\text{(or $O(\log N)$ with ANN indexes like HNSW)}.
\]

\subsection{Measured Extrapolation}

Using the measured per-pair Cross-Encoder time of $\sim$4.9\,ms
(derived from 0.57\,s~$\div$~116~chunks), Table~\ref{tab:scaling}
extrapolates wall-clock times:

\begin{table}[H]
\centering
\caption{Scaling extrapolation: Exhaustive Cross-Encoder vs.\ Two-Stage RAG
         (from Cell~12 output).}
\label{tab:scaling}
\begin{tabular}{@{}lrllr@{}}
\toprule
\textbf{Corpus Scale} & \textbf{$N$} & \textbf{Exhaustive CE} & \textbf{Two-Stage RAG} & \textbf{Speedup} \\
\midrule
Current IFAB subcorpus     & 116         & 0.57\,s   & 54.3\,ms & $10\times$ \\
Enterprise documentation   & 10\,000     & 49.31\,s  & 54.3\,ms & $907\times$ \\
Institutional archive      & 100\,000    & 8.2\,min  & 54.3\,ms & $9{,}079\times$ \\
Web-scale corpus           & 1\,000\,000 & 1.4\,hrs  & 54.3\,ms & $90{,}793\times$ \\
\bottomrule
\end{tabular}
\end{table}

\subsection{Conclusion}

Exhaustive Cross-Encoder search is infeasible beyond small corpora: at
100\,000 chunks it already requires \textbf{8.2 minutes} per query, and at
web scale (1M chunks) it requires \textbf{1.4 hours}.  The Two-Stage RAG
architecture achieves the best of both worlds:

\begin{enumerate}[nosep]
    \item \textbf{Recall via Bi-Encoder:} sub-millisecond approximate
          candidate filtering over the full collection.
    \item \textbf{Precision via Cross-Encoder:} high-fidelity re-ranking of a
          small, fixed-size candidate pool ($K{=}10$), keeping re-ranking cost
          constant at $\sim$36\,ms regardless of collection size.
\end{enumerate}

This decouples recall (cheap, scalable) from precision (expensive, bounded),
keeping total query latency under 55\,ms regardless of collection size.


% ───────────────────────────────────────────────────────────
\section{Extensions and Ablation Studies}
\label{sec:extensions}

\subsection{Candidate Pool Depth Ablation ($K \in \{5, 10, 20, 50\}$)}

We evaluated second-stage sensitivity to candidate pool depth $K \in \{5, 10, 20, 50\}$ across all 20 benchmark queries (Cell~13).

\begin{table}[H]
\centering
\small
\caption{Candidate pool depth ablation (from Cell~13 output).}
\label{tab:ablation}
\begin{tabular}{@{}cccc@{}}
\toprule
\textbf{Pool Size ($K$)} & \textbf{Mean Recall@3} & \textbf{Mean Latency (ms)} & \textbf{P95 Latency (ms)} \\
\midrule
5  & 0.550 & 36.31 & 45.72 \\
10 & 0.700 & 45.16 & 47.90 \\
20 & 0.700 & 79.11 & 81.34 \\
50 & 0.700 & 167.95 & 182.99 \\
\bottomrule
\end{tabular}
\end{table}

\noindent
Expanding $K$ from 5 to 10 yields a $+0.15$ gain in Recall@3 ($0.550 \to 0.700$). However, expanding to $K=20$ or $K=50$ produces \emph{zero additional gain}, as the remaining 30\% of gold chunks are absent from the top-50 pool entirely. Because Cross-Encoder inference grows linearly with $K$ ($\approx 3.6$\,ms per candidate pair, reaching 167.95\,ms at $K=50$), $K=10$ represents the empirical Pareto-optimal operating point.

\subsection{Hybrid Retrieval (BM25 + Dense RRF)}

We evaluated sparse lexical matching (BM25Okapi) fused with dense BGE vectors via Reciprocal Rank Fusion (Cell~14), defined as $\text{RRF}(d) = \frac{1}{60 + r_{\text{dense}}(d)} + \frac{1}{60 + r_{\text{bm25}}(d)}$, followed by Cross-Encoder re-ranking. On Keyword Trap queries, both Dense+CE and Hybrid+CE achieved $\text{Hit@3} = 1.00$ (4/4 queries). While dense re-ranking proved sufficient here, hybrid fusion offers a vital safeguard against lexical dropout when querying exact statutory numbers or domain jargon.

% ───────────────────────────────────────────────────────────
\section*{Code and Reproducibility}
\addcontentsline{toc}{section}{Code and Reproducibility}

All code, evaluation datasets, pre-indexed vectors, and the interactive execution notebook are available on GitHub:
\begin{center}
    \url{https://github.com/shobhan-iiit/Rag1}
\end{center}
The repository includes the complete standalone Jupyter notebook (\texttt{rag\_reranking.ipynb}), modular source code in \texttt{src/}, evaluation queries in \texttt{queries.json}, and pre-computed artifacts for offline verification.

\end{document}
